\section{结论与讨论}\label{ux7ed3ux8bbaux4e0eux8ba8ux8bba}

\section*{Conclusion and Discussion}

本文提出了一种面向多标签眼科疾病识别的可解释与隐私保护单域泛化框架，针对临床多疾病共存、多机构数据难以联合训练且跨机构部署存在域偏移、以及隐私泄露风险等问题，构建了一种统一解决框架。核心贡献包括：（1）自适应标签-特征关联性构建模块（ARC），利用Transformer自注意力机制和交叉注意力机制捕捉标签-标签关联性与标签-特征关联性，并通过注意力一致性损失对齐双分支输出，有助于减少多标签识别中的误检与漏检；（2）双分支引导域不变特征提取模块（DFE），通过风格增强模拟潜在域偏移、CAM提供疾病区域先验引导，并结合空间/通道注意力与 mask 过滤提取域不变特征，在提升未知域泛化性能的同时增强模型可解释性；（3）自适应高斯差分隐私机制（AdaGDP），通过动态衰减噪声强度，在中等隐私预算（\(\epsilon \approx 10\)）下控制性能损失，实现隐私保护与模型性能的平衡。系统性的实验验证了所提方法的有效性。尽管取得一定进展，本框架仍存在局限：类别不平衡可能影响少数类疾病识别的稳定性，未来可结合重采样或损失加权进一步缓解；当前跨域泛化主要基于共有标签空间展开，未来可探索开放类别和更复杂分布外场景下的泛化能力。

We propose a Privacy-Preserving Interpretable Single Domain Generalization framework for multi-label ocular disease recognition. The framework provides a unified solution to clinical multi-disease co-occurrence, restricted joint training with multi-institutional data, domain shift in cross-institutional deployment, and privacy leakage risk. The core contributions are as follows: (1) the Adaptive Label-Feature Relevance Construction (ARC) module uses Transformer self-attention and cross-attention mechanisms to capture label-label relevance and label-feature relevance, and aligns dual-branch outputs through attention consistency loss, helping reduce false positives and false negatives in multi-label recognition; (2) the Dual-Branch Guided Domain-Invariant Feature Extraction (DFE) module simulates potential domain shift through style augmentation, uses CAM to provide disease-region prior guidance, and combines spatial/channel attention with mask filtering to extract domain-invariant features, improving unseen-domain generalization performance while enhancing model interpretability; and (3) the Adaptive Gaussian Differential Privacy (AdaGDP) mechanism dynamically decays noise intensity to control performance loss under a moderate privacy budget (\(\epsilon \approx 10\)), achieving a balance between privacy preservation and model performance. Systematic experiments validate the effectiveness of the proposed method. Despite these advances, the framework still has limitations. Class imbalance may affect the stability of recognition for minority disease classes, which could be further alleviated by resampling or loss weighting in future work. In addition, current cross-domain generalization is mainly based on shared label spaces, and future work can explore generalization under open-category and more complex out-of-distribution scenarios.
